\subsection{LIE ALGEBRA OF A LIE GROUP}

Let \(\mathbf{G}\) be a Lie group. In \(\mathrm{U}(\mathbf{G})\) , as in any associative algebra, we write \([t, t'] = t * t' - t' * t\) . As \(\mathrm{T}_{e}(\mathbf{G})\) is the set of primitive elements of \(\mathrm{U}(\mathbf{G})\) , \([\mathrm{T}_{e}(\mathbf{G}), \mathrm{T}_{e}(\mathbf{G})] \subset \mathrm{T}_{e}(\mathbf{G})\) (Chapter II, § 1, no. 2, Proposition 4). The restriction of the bracket to \(\mathrm{T}_{e}(\mathbf{G})\) therefore defines on \(\mathrm{T}_{e}(\mathbf{G})\) a Lie algebra structure.

Lemma 1. Let \(\mathbf{X}\) and \(\mathbf{X}'\) be complete normable spaces, \(\mathbf{X}_0\) an open neighbourhood of 0 in \(\mathbf{X}\) and \(f\) an analytic mapping of \(\mathbf{X}_0\) into \(\mathbf{X}'\) such that \(f(0) = 0\) . Let \(f = f_1 + f_2 + f_3 + \dots\) be the expansion of \(f\) as an integral series about 0, where \(f_i\) is a homogeneous continuous-polynomial of degree \(i\) on \(\mathbf{X}\) with values in \(\mathbf{X}'\) . Let \(t\) be an element of \(\mathrm{TS}^2 (\mathbf{X})\) , considered as a point distribution on \(\mathbf{X}\) with support contained in \(\{0\}\) . Let \(t' = f_{\star}(t) \in \mathrm{TS}(\mathbf{X}')\) . The homogeneous component of \(t'\) of degree 1 is \(\langle f_2, t \rangle\) .

Let \(t_{1}^{\prime}\) be this component. Then, for every continuous linear mapping u of \(X^{\prime}\) into a polynormed space,

\[
\begin{array}{l l} u (t _ {1} ^ {\prime}) = \langle t ^ {\prime}, u \rangle & \text { because   } u \text {   is   continuous   and   linear } \\ = \langle t, u \circ f \rangle & (\text { Diff.   \&   Anal.   Man.,   R,   13.2.3 }) \\ = \langle t, u \circ f _ {2} \rangle & \text { because   } t \in \mathsf {T S} ^ {2} (\mathbf {X}) \\ = u (\langle t, f _ {2} \rangle) & (\text { Diff.   \&   Anal.   Man.,   R,   13.2.2 }), \end{array}
\]

whence the lemma.

PROPOSITION 24. Let G be a Lie group and (U, \(\phi\) , E) a chart on G such that \(\phi(e) = 0\) . Let V be an open neighbourhood of \(e\) such that \(\mathrm{V}^2 \subset \mathrm{U}\) . Let \(m\) be the analytic mapping \((a, b) \mapsto \phi(\phi^{-1}(a)\phi^{-1}(b))\) of \(\phi(V) \times \phi(V)\) into E. Let

\[
m = \sum_ {i, j \geqslant 0} m _ {i, j}
\]

be the expansion of \(m\) as an integral series about \((0,0)\) , where \(m_{i,j}\) is a bihomogeneous continuous-polynomial of bidegree \((i,j)\) on \(\mathbf{E} \times \mathbf{E}\) with values in \(\mathbf{E}\) .

\begin{enumerate}
\def\labelenumi{(\roman{enumi})}
\item
  \(m_{i,0} = m_{0,j} = 0\) for all \(i \neq 1\) and \(j \neq 1\) .
\item
  \(m_{1,0}(a,b)=a\) and \(m_{0,1}(a,b)=b\) for all \(a\in E\) , \(b\in E\) .
\item
  Let \(\psi: \mathrm{T}_e(\mathrm{G}) \to \mathrm{E}\) be the differential of \(\phi\) at \(e\) . For all \(u, v\) in \(\mathrm{T}_e(\mathrm{G})\) ,
\end{enumerate}

\[
\psi ([ u, v ]) = m _ {1, 1} (\psi (u), \psi (v)) - m _ {1, 1} (\psi (v), \psi (u)).
\]

\(m(a,0) = a\) , \(m(0,b) = b\) for all \(a\) , \(b\) in \(\phi (\mathrm{V})\) , which proves (i) and (ii). Let \(u, v\) be in \(\mathrm{T_e(G)}\) . Let \(\mathrm{T_0(E)}\) be identified with E and hence \(\psi\) with \(\mathrm{T_e}(\phi)\) . The images of \(u\) and \(v\) under \(\mathrm{T_e}(\phi)\) are \(\psi (u)\) and \(\psi (v)\) . The tensor product point distribution of these images is the symmetric product of \((\psi(u),0)\) and \((0,\psi (v))\) in \(\mathsf{TS(E\times E)} = \mathsf{TS(E)}\times \mathsf{TS(E)}\) , that is

\[
(\psi (u), 0) \otimes (0, \psi (v)) + (0, \psi (v)) \otimes (\psi (u), 0).
\]

Hence \(\phi_{*} (u * v)\) is the image of the above element under the mapping \(m\) of \(\phi(V) \times \phi(V)\) into \(E\) . Its component of degree 1 in TS(E) is, by Lemma 1,

\[
x = \langle m _ {1, 1}, (\psi (u), 0) \otimes (0, \psi (v)) + (0, \psi (v)) \otimes (\psi (u), 0) \rangle .
\]

We define a bilinear mapping \(n\colon (\mathbf{E}\times \mathbf{E})^2\to \mathbf{E}\) by

\[
n ((a, b), (a ^ {\prime}, b ^ {\prime})) = m _ {1, 1} (a, b ^ {\prime}).
\]

Then \(n((a, b), (a, b)) = m_{1,1}(a, b)\) and hence

\[
\begin{array}{l} x = \langle n, (\psi (u), 0) \otimes (0, \psi (v)) + (0, \psi (v)) \otimes (\psi (u), 0) \rangle \\ \quad = m _ {1, 1} (\psi (u), \psi (v)) + m _ {1, 1} (0, 0) = m _ {1, 1} (\psi (u), \psi (v)). \end{array}
\]

Similarly, \(\phi_{*} (v * u)\) admits \(m_{1,1}(\phi(v), \phi(u))\) as component of degree 1 in TS(E). As \(\phi([u, v])\) is of degree 1, this proves (iii).

COROLLARY. The normable space \(\mathrm{T}_e(\mathrm{G})\) together with the bracket, is a normable Lie algebra.

DEFINITION 6. The normable space \(\mathrm{T_e(G)}\) , together with the bracket, is called the normable Lie algebra of \(\mathbf{G}\) , or simply the Lie algebra of \(\mathbf{G}\) , and is denoted by \(\mathbf{L(G)}\) .

PROPOSITION 25. Let \(\mathbf{G}\) be a Lie group and \(\operatorname{E}(\mathbf{G})\) the enveloping algebra of \(\mathbf{L}(\mathbf{G})\) . The canonical injection of \(\mathbf{L}(\mathbf{G})\) into \(\operatorname{E}(\mathbf{G})\) defines a homomorphism \(\theta\) of the algebra \(\operatorname{E}(\mathbf{G})\) into the algebra \(\operatorname{U}(\mathbf{G})\) . If \(\mathbf{K}\) is of characteristic \(0\) , \(\eta\) is a bigebra isomorphism.

The bigebra U(G) is cocommutative (Differentiable and Analytic Manifolds, R, 13.5.1) and the filtration (U \(_{s}\) (G)) is compatible with the bigebra structure. The set of primitive elements of U(G) is L(G). It then suffices to apply chapter II, § 1, no. 6, Theorem 1.

When \(\mathbf{K}\) is of characteristic 0 we shall in future identify \(\mathrm{U}(\mathrm{G})\) with the enveloping algebra of L(G). By (2) and Proposition 7 (ii), the mapping \(t \mapsto t^{\vee}\) of U(G) into U(G) is then identified with the principal antiautomorphism of U(G) (Chapter I, § 2, no. 4).

PROPOSITION 26. Suppose that K is of characteristic \(p > 0\) . For all \(a \in \mathrm{L}(\mathrm{G})\) , \(a^p \in \mathrm{L}(\mathrm{G})\) and \(\operatorname{ad}(a^p) = (\operatorname{ad} a)^p\) (the power \(a^p\) being calculated in \(\mathrm{U}(\mathrm{G})\) ).

If \(a \in \mathrm{L}(\mathrm{G})\) , \(a\) is primitive in \(\mathrm{U}(\mathrm{G})\) , hence \(a^p\) is primitive in \(\mathrm{U}(\mathrm{G})\) (Chapter II, § 1, no. 2, Remark 1) and hence \(a^p \in \mathrm{L}(\mathrm{G})\) . Let \(\sigma_a\) (resp. \(\tau_a\) ) be the linear mapping \(x \mapsto a * x\) (resp. \(x \mapsto x * a\) ) of \(\mathrm{U}(\mathrm{G})\) into \(\mathrm{U}(\mathrm{G})\) . For all \(x \in \mathrm{U}(\mathrm{G})\) , (ad \(a\) ) \((x) = (\sigma_a - \tau_a)(x)\) and hence (ad \(a\) ) \(^p = (\sigma_a - \tau_a)^p\) . But \(\sigma_a\) and \(\tau_a\) commute and therefore \((\sigma_a - \tau_a)^p = (\sigma_a)^p - (\tau_a)^p = \tau_a^p - \sigma_a^p\) , whence the second assertion.

DEFINITION 7. Let X be a manifold of class \(\mathbf{C}^r\) ( \(r \geqslant 2\) ) and \(\mathfrak{g}\) a complete normable Lie algebra. A law of infinitesimal left (resp. right) operation of class \(\mathbf{C}^{r-1}\) of \(\mathfrak{g}\) on X is a mapping \(a \mapsto \mathbf{D}_a\) of \(\mathfrak{g}\) into the set of vector fields on X with the following properties:

\begin{enumerate}
\def\labelenumi{(\alph{enumi})}
\item
  the mapping \((a, x) \mapsto \mathrm{D}_a(x)\) is a morphism of class \(\mathbf{C}^{r-1}\) of the trivial vector bundle \(g \times X\) into the vector bundle \(\mathrm{T}(X)\) ;
\item
  \([\mathrm{D}_a,\mathrm{D}_b] = -\mathrm{D}_{[a,b]}\) (resp. \([\mathrm{D}_a,\mathrm{D}_b] = \mathrm{D}_{[a,b]})\) for all \(a,b\) in \(\mathfrak{g}\) .
\end{enumerate}

In particular, each vector field \(\mathbf{D}_a\) is of class \(\mathbf{C}^{r - 1}\) .

Remark. Let X be a manifold of class \(C^{r}\) , g a finite-dimensional Lie algebra and \(a \mapsto D_{a}\) a linear mapping of g into the vector space of vector fields of class \(C^{r-1}\) on X. Then condition (a) of Definition 7 holds. For, by considering a basis of g and applying Differentiable and Analytic Manifolds, R, 7.7.1, the problem is reduced to the case dim g = 1 and our assertion is then obvious.

PROPOSITION 27. Let G be a Lie group and X a manifold of class \(\mathbf{C}^r\) . Suppose that a law of left (resp. right) operation of class \(\mathbf{C}^r\) of G on X is given. For all \(a \in \mathrm{L}(\mathrm{G})\) , let \(\mathrm{D}_a\) be the field of point distributions defined by \(a\) on X.

\begin{enumerate}
\def\labelenumi{(\roman{enumi})}
\item
  The mapping \((a, x) \mapsto \mathrm{D}_a(x)\) is a morphism of class \(\mathbf{C}^{r-1}\) of the trivial vector bundle \(\mathbf{L}(\mathbf{G}) \times \mathbf{X}\) into the vector bundle \(\mathrm{T}(\mathbf{X})\) .
\item
  Let \(\mathbf{I}\) be an open subset of \(\mathbf{K}\) containing 0 and \(\gamma: \mathbf{I} \to \mathbf{G}\) a mapping of class \(\mathbf{C}^r\) such that \(\gamma(0) = e\) . Let \(a = \mathrm{T}_0(\gamma) 1 \in \mathbf{L}(\mathbf{G})\) . If \(f\) is a function of class \(\mathbf{C}^r\) on an open subset of \(\mathbf{X}\) , then
\end{enumerate}

\((D_{a}f)(x) = \lim_{k\in \mathbb{K}^{*},k\to 0}k^{-1}(f(\gamma (k)x) - f(x))\) if G operates on the left,

\[
(D _ {a} f) (x) = \lim _ {k \in K ^ {*}, k \rightarrow 0} k ^ {- 1} (f (x \gamma (k)) - f (x)) \quad \text {if } G \text { operates on the right}.
\]

\begin{enumerate}
\def\labelenumi{(\roman{enumi})}
\setcounter{enumi}{2}
\tightlist
\item
  If \(r \geqslant 2\) , the mapping \(a \mapsto \mathbf{D}_a\) is a law of infinitesimal left (resp. right) operation of class \(\mathbf{C}^{r-1}\) of \(\mathbf{L}(\mathbf{G})\) on \(\mathbf{X}\) .
\end{enumerate}

Suppose that G operates on X on the left. Let \(\phi: G \times X \to X\) be the law of operation. Then \(T(\phi)\) is a \(\phi\) -morphism of class \(C^{r-1}\) of the vector bundle \(T(G) \times T(X)\) into the vector bundle \(T(X)\) (Differentiable and Analytic

Manifolds, R, 8.1.2). The induced vector bundle \((\mathrm{T}(\mathrm{G}) \times \mathrm{T}(\mathrm{X}))|(\{e\} \times \mathrm{X})\) is identified with \(\mathrm{E} = \mathrm{L}(\mathrm{G}) \times \mathrm{T}(\mathrm{X})\) . Hence \(\mathrm{T}(\phi)|\mathrm{E}\) is a vector bundle morphism of class \(\mathbf{C}^{r-1}\) . For \((a, x) \in \mathrm{L}(\mathrm{G}) \times \mathrm{X}\) , \(\mathrm{T}(\phi)(a, x) = \mathbf{D}_{a}(x)\) , whence (i).

The formula giving \((\mathrm{D}_a f)(x)\) follows from §2, the end of no. 2, and Differentiable and Analytic Manifolds, R, 8.4.5.

Suppose that \(r \geqslant 2\) . Let \(a, b\) be in \(\mathbf{L}(\mathbf{G})\) and \(f\) be a function of class \(\mathbf{C}^r\) on an open subset of \(\mathbf{X}\) . Then

\[
\begin{array}{l l} \mathrm{D} _ {[ a, b ]} f = \mathrm{D} _ {b} (\mathrm{D} _ {a} f) - \mathrm{D} _ {a} (\mathrm{D} _ {b} f) & \text { by   (17) } \\ = [ \mathrm{D} _ {b}, \mathrm{D} _ {a} ] f & (\text { Diff.   \&   Anal.   Man.,   R,   8.5.3 }). \end{array}
\]

Let \(x \in X\) . By taking \(f\) to be a chart on an open neighbourhood of \(x\) , it follows that \(D_{[a,b]}(x) = [D_b, D_a](x)\) , whence (iii). The argument is similar if \(G\) operates on \(X\) on the right.

When \(r \geqslant 2\) , the mapping \(a \mapsto D_{a}\) is called the law of infinitesimal operation associated with the given law of operation.

